\subsection*{ソフトウェア構成変更統制}
\addcontentsline{toc}{subsection}{3 ソフトウェア構成変更統制}
ソフトウェア構成変更統制は、ライフサイクル中に構成品目へ必要となる変更を扱う活動である。何を変更するか、誰が承認するか、どのように実施するか、要求からの逸脱や免除をどう扱うかを定める。

変更統制から得られる情報は、変更量、破壊的変更、再作業、欠陥修正、拡張要求の測定にも使える。

\subsubsection*{ソフトウェア変更の要求・評価・承認}
\addcontentsline{toc}{subsubsection}{3.1 ソフトウェア変更の要求・評価・承認}
管理対象品目への変更は、通常、変更要求として提出される。ソフトウェア変更要求プロセスは、変更要求の提出、記録、費用と影響の評価、承認・修正・延期・却下を扱う正式な手順である。

\term{変更要求}{Change Request}は、ライフサイクルのどの時点でも、誰からでも発生し得る。問題報告に対する是正処置、機能追加、範囲変更、スケジュール変更、費用変更、手順変更などが原因になり得る。

変更要求を受け取ると、技術評価または影響分析を行う。どの構成品目が影響を受けるか、どの設計・コード・テスト・文書を変更する必要があるかを調べる。ここで、構成品目間の関係情報が重要になる。

\par\smallskip
\begin{center}
\centering
\IfFileExists{assets/generated_figures/ch08/fig-ch08-08.png}{\includegraphics[width=0.96\linewidth]{assets/generated_figures/ch08/fig-ch08-08.png}}{\fbox{\parbox[c][48mm][c]{0.9\linewidth}{\centering 画像生成待ち\\変更要求プロセスの流れ}}}
\par\smallskip
\refstepcounter{figure}\label{fig:ch08-08}図 \thefigure: 変更要求プロセスの流れ
\end{center}
図 \thefigure は、変更要求プロセスの流れを視覚的に整理したものである。
\par\smallskip
\paragraph*{ソフトウェア構成管理委員会}
\addcontentsline{toc}{subsubsection}{3.1.1 ソフトウェア構成管理委員会}
変更の承認または却下の権限は、構成管理委員会に置かれる。小規模プロジェクトでは、委員会ではなくプロジェクトリーダーや指定された個人が権限を持つこともある。変更の重要性、予算・日程への影響、ライフサイクル上の時点によって、承認権限の階層を分ける場合がある。

委員会の範囲がソフトウェアに限られる場合は、ソフトウェア構成管理委員会と呼ぶ。SCM 代表者は常に関与し、必要な利害関係者が委員会レベルに応じて参加する。

\paragraph*{ソフトウェア変更要求プロセス}
\addcontentsline{toc}{subsubsection}{3.1.2 ソフトウェア変更要求プロセス}
有効な SCR プロセスには、変更要求の発行、ワークフローの強制、委員会判断の記録、変更処理情報の報告を支援する手順とツールが必要である。問題報告システムと連携すると、問題解決の追跡と解決速度の把握がしやすい。

\paragraph*{ソフトウェア変更要求フォーム}
\addcontentsline{toc}{subsubsection}{3.1.3 ソフトウェア変更要求フォーム}
変更要求には、要求の内容と根拠を示すフォームが必要である。承認された場合には、変更の認証や完了確認に使うフォームも必要になる。ツールで管理する場合でも、何を入力し、誰が責任を持つかは人が設計する。

\subsubsection*{ソフトウェア変更の実施}
\addcontentsline{toc}{subsubsection}{3.2 ソフトウェア変更の実施}
承認された変更要求は、定義済みの手順とスケジュールに従って実施する。複数の変更が同時に実施される場合、どの変更要求がどの版や基準版に含まれたかを追跡できなければならない。

変更完了時には、構成監査や品質検証を行い、承認済み変更だけが行われたことを確認する。版管理ツールは、チェックアウト、変更、記録、保存、差し戻しを支援する。並行開発や分散開発では、より強力なツールが必要になる。

\subsubsection*{逸脱と免除}
\addcontentsline{toc}{subsubsection}{3.3 逸脱と免除}
逸脱は、品目を作る前に、特定の性能要求または設計要求から一時的に外れることを認める書面上の許可である。免除は、生産中または検査後に見つかった問題に対して、指定要求から外れていても使用に適すると判断して受け入れる書面上の許可である。

\par\smallskip\noindent\refstepcounter{table}\label{tab:ch08-05}表 \thetable: 逸脱と免除における区分・時点・意味\par\smallskip
表 \thetable は、逸脱と免除における区分・時点・意味を比較・確認しやすい形で整理したものである。\par\smallskip
\begin{longtable}{p{0.20\linewidth}p{0.25\linewidth}p{0.49\linewidth}}\toprule
区分 & 時点 & 意味 \\ \midrule
\endfirsthead\toprule
区分 & 時点 & 意味 \\ \midrule
\endhead
逸脱 & 品目を作る前 & 要求を満たせないことが事前にわかり、一定範囲で外れることを認める。 \\
免除 & 生産中または検査後 & 要求から外れた品目を、そのまままたは承認済み手直し後に使用可とする。 \\
\bottomrule\end{longtable}
